% пример введения

В условиях стремительного развития цифровых образовательных 
технологий и увеличения объёма информации, которую необходимо усваивать 
современным студентам, традиционные подходы к обучению демонстрируют 
недостаточную эффективность. Интеллектуальные обучающие системы (ИОС) 
представляют собой перспективное направление, позволяющее персонализировать 
образовательный процесс, адаптировать подачу материала под индивидуальные 
особенности обучающегося и повысить мотивацию к обучению.

Особую актуальность приобретает разработка конструкторов адаптивных курсов, 
которые учитывают не только когнитивные особенности студентов (стиль обучения, 
уровень подготовки, психологические характеристики), но и контекстные данные 
(географическое положение, погодные условия, время суток), влияющие на способность
к усвоению материала. Такой комплексный подход позволяет создать по-настоящему 
персонализированную образовательную среду.

Объект исследования – процесс разработки интеллектуальной обучающей системы с 
использованием фреймворка Django.Предмет исследования – модели и методы 
построения адаптивных образовательных траекторий на основе когнитивного профиля 
обучающегося и контекстных данных.Цель работы – провести анализ существующих 
моделей обучения и разработать оригинальную модель для интеллектуальной 
обучающей системы, учитывающую когнитивный профиль студента и внешние контекстные 
факторы. 